% !TeX root = ../eccv2016submission.tex

Many attempts have been made to improve the training of very deep networks. Earlier works adopted greedy layer-wise training or better initialization schemes
to alleviate the vanishing gradients and diminishing feature reuse problems \cite{fahlman1989cascade,erhan2010does,glorot2010understanding}.
A notable recent contribution towards training of very deep networks is Batch Normalization~\cite{ioffe2015batch}, which standardizes the mean and variance of hidden layers with respect to each mini-batch. This approach reduces the vanishing gradients problem and yields a strong regularizing effect.

Recently, several authors introduced extra skip connections to improve the information flow during forward and backward propagation. Highway Networks \cite{srivastava2015highway} allow earlier representations to flow unimpededly to later layers through parameterized skip connections known as ``information highways'', which can cross several layers at once. The skip connection parameters, learned during training, control the amount of information allowed on these ``highways''.

\paragraph{\textbf{Residual networks (ResNets)}\cite{he2015deep}} simplify Highway Networks by shortcutting (mostly) with identity functions. This simplification greatly improves training efficiency, and enables more direct feature reuse. ResNets are motivated by the observation that neural networks tend to obtain \emph{higher training error} as the depth increases to very large values. This is counterintuitive, as the network gains more parameters and therefore better function approximation capabilities. The authors conjecture that the networks become \emph{worse} at function approximation because the gradients and training signals vanish when they are propagated through many layers. As a fix, they propose to add \emph{skip connections} to the network. Formally, if $H_\ell$ denotes the output of the $\ell^{th}$ layer (or sequence of layers) and $f_\ell(\cdot)$ represents a\ typical convolutional transformation from layer $\ell\!-\!1$ to $\ell$, we obtain
\begin{equation}
\label{eq:res_block}
 H_{\ell} = {\tt ReLU} (f_\ell( H_{\ell-1}) + \text{id} (H_{\ell-1})),
\end{equation}
where $\text{id}(\cdot)$ denotes the identity transformation and we assume a ${\tt ReLU}$ transition function~\cite{nair2010rectified}. Fig.~\ref{fig:resblock} illustrates an example of a function $f_\ell$, which consists of multiple convolutional and Batch Normalization layers.
When the output dimensions of $f_\ell$ do not match those of $H_{\ell-1}$, the authors redefine $\text{id}(\cdot)$ as a linear projection to reduce the dimensions of $\text{id}(H_{\ell-1})$ to match those of $f_\ell(H_{\ell-1})$.
The propagation rule in \eqref{eq:res_block} allows the network to pass gradients and features (from the input or those learned in earlier layers) back and forth between the layers via the identity transformation $\text{id}(\cdot)$.

\paragraph{\textbf{Dropout.}} Stochastically dropping hidden nodes or connections has been a popular regularization method for neural networks. The most notable example is Dropout \cite{srivastava2014dropout}, which multiplies each hidden activation by an independent Bernoulli random variable. Intuitively, Dropout reduces the effect known as ``co-adaptation'' of hidden nodes collaborating in groups instead of independently producing useful features; it also makes an analogy with training an ensemble of exponentially many small networks. Many follow up works have been empirically successful, such as DropConnect \cite{icml2013_wan13}, Maxout \cite{goodfellow2013maxout} and DropIn \cite{smith2015gradual}.

Similar to Dropout, \name{} can be interpreted as training an ensemble of networks, but with different depths, possibly achieving higher diversity among ensemble members than ensembling those with the same depth. Different from Dropout, we make the network shorter instead of thinner, and are motivated by a different problem. Anecdotally, Dropout loses effectiveness when used in combination with Batch Normalization~\cite{ioffe2015batch,blogpostcifar10}. Our own experiments with various Dropout rates (on CIFAR-10) show that Dropout gives practically no improvement when used on 110-layer ResNets with Batch Normalization.

We view all of these previous approaches to be extremely valuable and consider our proposed training with \name{} complimentary to these efforts. In fact, in our experiments we show that training with \name{} is indeed very effective on ResNets with Batch Normalization.
